\documentclass{article}

\usepackage{amsmath,amssymb}
\usepackage{xurl}
\usepackage[hidelinks]{hyperref}

% Shared commands for notes in docs/paper-gaps/.

\DeclareUnicodeCharacter{2080}{\ifmmode{}_{0}\else\textsubscript{0}\fi}
\DeclareUnicodeCharacter{2081}{\ifmmode{}_{1}\else\textsubscript{1}\fi}
\DeclareUnicodeCharacter{2082}{\ifmmode{}_{2}\else\textsubscript{2}\fi}
\DeclareUnicodeCharacter{2099}{\ifmmode{}_{n}\else\textsubscript{n}\fi}

\newcommand{\Lean}{\textsc{Lean}}
\ExplSyntaxOn
\NewDocumentCommand{\leanid}{m}{
  \ifmmode
    \textsf{\detokenize{#1}}
  \else
    \group_begin:
    \urlstyle{sf}
    \tl_set:Nn \l_tmpa_tl {#1}
    \tl_replace_all:Nnn \l_tmpa_tl {\_} {_}
    \exp_args:NV \path \l_tmpa_tl
    \group_end:
  \fi
}
\ExplSyntaxOff
\providecommand{\tr}{\operatorname{tr}}
\newcommand{\tnleanIssueBase}{https://github.com/LionSR/TNLean/issues/}
\newcommand{\tnleanPrBase}{https://github.com/LionSR/TNLean/pull/}
\newcommand{\tnleanDocBase}{https://sirui-lu.com/TNLean/docs/}
\newcommand{\ghissue}[1]{\href{\tnleanIssueBase#1}{GitHub issue \##1}}
\newcommand{\ghpr}[1]{\href{\tnleanPrBase#1}{PR \##1}}
\newcommand{\doclink}[1]{\href{\tnleanDocBase#1.html}{\path{#1}}}

% Dirac notation and the identity operator, as in blueprint/src/macros/common.tex.
% \providecommand keeps note-local definitions authoritative.
\usepackage{dsfont}
\providecommand{\ket}[1]{|#1\rangle}
\providecommand{\bra}[1]{\langle #1|}
\providecommand{\braket}[2]{\langle #1 | #2 \rangle}
\providecommand{\ketbra}[2]{|#1\rangle\!\langle #2|}
\providecommand{\Id}{\mathds{1}}
\providecommand{\one}{\mathds{1}}
\providecommand{\C}{\mathbb{C}}
\providecommand{\MN}[1]{M_{#1}(\C)}
\providecommand{\MD}{M_D(\C)}

% External referencing for paper sources.  The build helper writes sanitized
% auxiliary files under build/paper-aux/ and excludes duplicated source labels.
\usepackage{xr}

% Import a generated paper auxiliary file with a caller-supplied label prefix.
% The candidate paths support builds from docs/paper-gaps/, the repository root,
% and build/paper-aux/ itself.
\newcommand{\importpaperaux}[2]{%
  \IfFileExists{../../build/paper-aux/#2.aux}{%
    \externaldocument[#1]{../../build/paper-aux/#2}%
  }{%
    \IfFileExists{build/paper-aux/#2.aux}{%
      \externaldocument[#1]{build/paper-aux/#2}%
    }{%
      \IfFileExists{#2.aux}{%
        \externaldocument[#1]{#2}%
      }{}%
    }%
  }%
}

% General external reference macros
\newcommand{\paperref}[2]{\ref{#1-#2}}
\newcommand{\papereqref}[2]{(\ref{#1-#2})}

% CPSV16 source (1606.00608 / MPDO-22-12-17-2.tex) external document and shortcuts
\importpaperaux{cpsv16-}{1606.00608/MPDO-22-12-17-2-tnlean}
\newcommand{\cpsvref}[1]{\paperref{cpsv16}{#1}}
\newcommand{\cpsveqref}[1]{\papereqref{cpsv16}{#1}}

% Verdict marker: \gapnote{<kind>}{<status>}. Kind is the policy
% classification (clarification, local-correction, scope-restriction,
% unfaithful, false-source, open-gap); status is open, wip (elimination
% underway), resolved, or historical. Machine-read by the site index;
% prints a verdict line.
\newcommand{\gapnote}[2]{\par\noindent\textbf{Verdict:} #1 (#2).\par}


\title{Zero-Weight Sector Completions in the MPDO Refinement Map}
\author{}
\date{2026-07-12}

\begin{document}
\maketitle
\gapnote{local-correction}{resolved}

\section{The undefined source quotient}

In the proof of Proposition~C.7 of
Ref.~\cite[Appendix~C.2, lines~1523--1535]{Cirac2016MPDO_arXiv}, the
refinement map on the sector pair $(k,h)$ is
\begin{align}
    \mathcal T_{k,h}(X)
        &=\frac{1}{a_kb_h}X\otimes
        \left[\bigoplus_l
        \left(\eta_{k,l}\otimes\eta_{l,h}\right)\right].
    \label{eq:cpgsv17_mpdo_zero_weight_preparation_completion_source_map}
\end{align}
Here $\eta_{k,h}$ is the neighboring-sector operator, while $a_k$ and $b_h$
are its scalar trace factors. Proposition~C.7 assumes
\begin{align}
    \eta_{k,h}
        &\succeq0,
    \label{eq:cpgsv17_mpdo_zero_weight_preparation_completion_eta_positive}\\
    \tr(\eta_{k,h})
        &=a_kb_h,
    \label{eq:cpgsv17_mpdo_zero_weight_preparation_completion_eta_trace}\\
    \sum_k a_kb_k
        &=1.
    \label{eq:cpgsv17_mpdo_zero_weight_preparation_completion_weight_normalization}
\end{align}
These hypotheses do not require every product $a_kb_h$ to be nonzero.
Consequently,
Eq.~\ref{eq:cpgsv17_mpdo_zero_weight_preparation_completion_source_map} is
undefined on a sector pair with $a_kb_h=0$.

\section{Trace-preserving completion}

Define the operator on the appended subspin space by
\begin{align}
    \Omega_{k,h}
        &=\bigoplus_l
        \left(\eta_{k,l}\otimes\eta_{l,h}\right).
    \label{eq:cpgsv17_mpdo_zero_weight_preparation_completion_omega_definition}
\end{align}
Positivity and the trace calculation in
Ref.~\cite{Cirac2016MPDO_arXiv} give
\begin{align}
    \Omega_{k,h}
        &\succeq0,
    \label{eq:cpgsv17_mpdo_zero_weight_preparation_completion_omega_positive}\\
    \tr(\Omega_{k,h})
        &=a_kb_h.
    \label{eq:cpgsv17_mpdo_zero_weight_preparation_completion_omega_trace}
\end{align}
If $a_kb_h=0$, then
Eqs.~\ref{eq:cpgsv17_mpdo_zero_weight_preparation_completion_omega_positive}
and~\ref{eq:cpgsv17_mpdo_zero_weight_preparation_completion_omega_trace}
imply
\begin{align}
    \Omega_{k,h}
        &=0.
    \label{eq:cpgsv17_mpdo_zero_weight_preparation_completion_zero_omega}
\end{align}

For an active sector pair, $a_kb_h>0$, append the density operator
$(a_kb_h)^{-1}\Omega_{k,h}$. For a zero-weight pair, append any density
operator on the same subspin space. These choices define a completely positive
trace-preserving map on every sector pair.

The arbitrary density operator on a zero-weight pair does not affect the
refinement identity. The preceding partial trace $\mathcal T_0$ produces the
factor
\begin{align}
    a_kb_hB_{k,h}(X)
        &=0
        \qquad\text{when }a_kb_h=0,
    \label{eq:cpgsv17_mpdo_zero_weight_preparation_completion_zero_partial_trace}
\end{align}
where $B_{k,h}(X)$ denotes the corresponding sector block. Equivalently, the
same block in the three-site closure is
\begin{align}
    B_{k,h}(X)\otimes\Omega_{k,h}
        &=0.
    \label{eq:cpgsv17_mpdo_zero_weight_preparation_completion_zero_three_site_block}
\end{align}
Thus every possible completion on an inactive pair gives the same
two-site-to-three-site refinement identity.

\section{Verdict}

The inactive-sector density operator supplies a local completion omitted from
the source formula. It adds no hypothesis to Proposition~C.7 of
Ref.~\cite{Cirac2016MPDO_arXiv} and does not change its two-site-to-three-site
identity. Formal declarations implementing this completion should carry a
\emph{Local fix (zero-weight quotient)} marker and cite this note.

\bibliographystyle{alpha}
\bibliography{references}

\end{document}
